\chapter{Metodologia}

A metodologia é a seção do trabalho que descreve os métodos e processos adotados para atingir os objetivos da pesquisa. Ela deve responder à pergunta: \textbf{“Como a pesquisa foi realizada?”}. Nesta seção, será explicado como compor uma metodologia utilizando a abordagem tradicional e, em seguida, como estruturá-la com o uso do \textbf{Design Science Research (DSR)}.

\section{Metodologia Tradicional}

Na abordagem tradicional, a metodologia é estruturada com base em uma **sequência lógica** que inclui:  

\subsection{Tipo de Pesquisa}  
Defina a natureza da pesquisa, que pode ser:  
\begin{itemize}
    \item \textbf{Exploratória}: Compreender um problema pouco estudado.  
    \item \textbf{Descritiva}: Detalhar características ou fenômenos.  
    \item \textbf{Experimental}: Testar hipóteses por meio de experimentos.  
    \item \textbf{Explicativa}: Identificar causas e efeitos.  
\end{itemize}

**Exemplo**:  
\textit{“Este trabalho caracteriza-se como uma pesquisa explicativa, com abordagem quantitativa, visando identificar padrões de disseminação do \textit{Aedes aegypti} em regiões afetadas.”}  

\subsection{Fontes e Coleta de Dados}  
Descreva as fontes dos dados e os métodos utilizados para coleta:  
\begin{itemize}
    \item **Dados Primários**: Entrevistas, experimentos, questionários, etc.  
    \item **Dados Secundários**: Bases de dados públicas, artigos científicos, etc.  
\end{itemize}

**Exemplo**:  
\textit{“Os dados climáticos foram coletados do INPE, enquanto os dados epidemiológicos foram obtidos do SINAN, entre 2010 e 2023.”}  

\subsection{Análise dos Dados}  
Explique as técnicas utilizadas para analisar os dados coletados:  
\begin{itemize}
    \item **Quantitativas**: Estatística descritiva, regressão, aprendizado de máquina.  
    \item **Qualitativas**: Análise de conteúdo, categorização de dados.  
\end{itemize}

**Exemplo**:  
\textit{“A análise dos dados foi realizada por meio de aprendizado de máquina, utilizando o algoritmo Random Forest com métricas de avaliação como acurácia e precisão.”}  

\subsection{Ferramentas Utilizadas}  
Liste os softwares, ferramentas e linguagens utilizadas.  

**Exemplo**:  
\begin{itemize}
    \item Software: Python 3.8  
    \item Bibliotecas: Pandas, Scikit-learn, Matplotlib  
    \item Ambiente: Jupyter Notebook  
\end{itemize}

\section{Metodologia Utilizando Design Science Research (DSR)}  
O \textbf{Design Science Research (DSR)} é uma abordagem voltada para pesquisas aplicadas que buscam criar e avaliar artefatos, como modelos, algoritmos ou sistemas. A metodologia DSR é estruturada em seis etapas principais:

\subsection{Identificação do Problema}  
Defina o problema prático ou teórico que será abordado.  

**Exemplo**:  
\textit{“O problema identificado é a falta de ferramentas preditivas para monitorar a distribuição de \textit{Aedes aegypti} em cenários de mudanças climáticas.”}  

\subsection{Objetivos da Solução}  
Estabeleça os objetivos que o artefato deve cumprir.  

**Exemplo**:  
\textit{“Desenvolver um modelo preditivo utilizando aprendizado de máquina para prever a distribuição do vetor com base em dados climáticos.”}  

\subsection{Desenvolvimento do Artefato}  
Descreva como a solução será desenvolvida, especificando técnicas e ferramentas.  

**Exemplo**:  
\textit{“O modelo será implementado com Random Forest, utilizando dados climáticos (INPE) e epidemiológicos (SINAN). O ambiente de desenvolvimento será Python 3.8.”}  

\subsection{Demonstração}  
Explique como a solução será testada em um cenário prático.  

**Exemplo**:  
\textit{“A demonstração será realizada com dados históricos do Pantanal, testando a capacidade do modelo em prever surtos de dengue em diferentes regiões.”}  

\subsection{Avaliação}  
Defina os critérios e métricas para avaliar a eficácia do artefato.  

**Exemplo**:  
\textit{“A avaliação será baseada em métricas como acurácia, sensibilidade e especificidade, comparando os resultados com modelos preexistentes.”}  

\subsection{Comunicação dos Resultados}  
Apresente como os resultados serão discutidos e documentados.  

**Exemplo**:  
\textit{“Os resultados serão apresentados por meio de gráficos comparativos, tabelas de métricas e discussões detalhadas sobre a eficácia do modelo.”}  

\section{Resumo Comparativo: Metodologia Tradicional vs. DSR}

\begin{tabular}{|p{6cm}|p{6cm}|}
\hline
\textbf{Metodologia Tradicional} & \textbf{Metodologia com DSR} \\ \hline
Define tipo de pesquisa e métodos tradicionais & Enfoca a criação e avaliação de um artefato \\ \hline
Coleta de dados e análise descritiva/quantitativa & Estrutura em etapas: problema, solução, avaliação \\ \hline
Foco na interpretação dos dados & Foco na solução prática e inovadora \\ \hline
Utiliza técnicas como estatísticas e entrevistas & Desenvolve e valida um artefato \\ \hline
\end{tabular}

A metodologia é um elemento essencial do trabalho acadêmico, devendo ser estruturada de forma lógica e detalhada. Enquanto a **abordagem tradicional** é amplamente utilizada em pesquisas teóricas e empíricas, o **Design Science Research (DSR)** é ideal para estudos aplicados, especialmente em áreas como computação, engenharia e tecnologia.


\section[Exemplo de capítulo com notação matemática]{Exemplo de capítulo com notação matemática}

Este capítulo contém alguns exemplos de ambientes matemáticos.

\section{Como criar uma definição}

\begin{defi}
Uma Topologia sobre o conjunto $X$ é uma família $\tau
\subset \beta(X)$ que satisfaz $ \mathbb{S} $:
\begin{enumerate}
 \item$\emptyset,X$ pertencem a $\tau;$
 \item Seja  $\{A_{\alpha} \in \tau /\alpha \in \Gamma\}$
 uma família arbitrária de subconjuntos de  $\tau$ então, a união dos elementos dessa família ainda pertence a  $\tau$, ou seja, $\bigcup \limits_{\alpha \in \Gamma} A_{\alpha}\in\tau;$
 \item Sejam $B_{1},B_{2},B_{3},..., B_{n} \in \tau$ uma família finita de subconjuntos de $\tau$, com $n \geq 1$, então
 interseção dessa família ainda pertence a $\tau$, ou seja, $\bigcap\limits_{i=1}^{n} B_{i}\in \tau.$
\end{enumerate}
\end{defi}

\section{Como criar uma observação}

\begin{obs}

\item Os elementos da Topologia  $\tau$ caracterizam os Conjuntos Abertos de $X.$

\item O par ordenado $(X,\tau)$ é chamado de Espaço Topológico.

\item Sim, de fato:

\item i) $\emptyset, X \in \tau$, pois são subconjuntos de $X;$

\item ii) Seja \{$J_{\lambda}\}_{\lambda\in L}$ temos que$J_{\lambda}\in \tau$ então $\bigcup\limits_{\lambda\in L} J_{\lambda}$ é
    um subconjunto de $X$ e portanto, pertence a $\tau;$

\item iii) Sejam $J_{1},J_{2},...J_{n} \in \tau$ temos que $J_{1}\bigcap J_{2}\bigcap ...J_{n}$ é um suconjunto de $X$ e portanto pertence a $\tau.$

\end{obs}

\section{Como criar um exemplo}

\begin{ex}{Topologia Euclidiana ou Usual} - $\tau_{us}$

Essa topologia é formada por intervalo do $\mathbb{R}^{n}$
Os dois tipos mais estudados são a Topologia Euclidiana Usual na Reta (ou  de $\mathbb{R}$) e a Topologia Euclidiana Usual (ou do  no Plano Complexo (ou do $\mathbb{R}^{2}$).

Seja a topologia $\tau = \{A \subset \mathbb{R}\}, \emptyset\}$ onde $X = \mathbb{R}$, dizemos que $A \in \tau$, se e somente se, para todo $y \in A$, existe um intervalo aberto $(c,d)$ tal que
$y \in (c,d) \subset A.$

\item i) É fácil notar que $\mathbb{R}, \emptyset \in \tau;$

\item ii) Consideremos a família $\{A_{\alpha} \in \tau / \alpha \in\tau\},$
queremos mostrar que, $\bigcup\limits_{\alpha\in \Gamma} A_{\alpha}\in \tau.$

De fato, considerando $y \in \bigcup\limits_{\alpha\in \Gamma} A_{\alpha}$ ,
desse modo, existe $\alpha_{0} \in \Gamma$ tal que $y \in A_{\alpha_{0}} \in \tau$, assim, existe $(c,d)$ e
temos que $y\in (c,d) \subset A_{\alpha_{0}}\subset \bigcup\limits_{\alpha\in \Gamma}A_{\alpha}.$

\item iii) Consideremos $B_{1},B_{2}$ subconjuntos finitos de $\tau$, desse modo
tomemos $y \in B_{1}\bigcap B_{2}$, assim $y\in B_{1}$ e $y\in B_{2}$,
logo existem $(c_{1}, d_{1})$ e $(c_{2}, d_{2})$ tais que $y$
$\in0 (c_{1}, d_{1}) \subset B_{1}$ e $y\in (c_{2}, d_{2}) \subset B_{2}.$

Vamos denotar por $c = max(c_{1},c_{2})$ e $d = min(d_{1},d_{2})$, assim temos:

\begin{center}
 $y \in (c,d) \subset B_{1} \bigcap B_{2}$
\end{center}
 
Por indução temos  se $B_{1},B_{2},..., B_{n}$ são conjuntos finitos de $\tau$, então, $\bigcap\limits_{i=1}^{n}B_{i}\in\tau.$
\end{ex}

\section{Como criar um teorema}

\begin{teo}: Seja $(X,\tau)$ um espaço topológico e $ \mathfrak{F}$ a
família de conjuntos fechados, então:
\begin{description}
  \item[i)]$\emptyset, X$ pertencem a $ \mathfrak{F}$;
  \item[ii)]Considere $F_{1}, F_{2},..., F_{n}$ conjuntos fechados
  em $X$; então a união finita $\bigcup\limits_{i=1}^{n} F_{i}$ é fechada
  em $X$;
  \item[iii)]Considere $\{ F_{\alpha} \in \mathfrak{F} / \alpha \in
  \Gamma\}$ uma família de elementos arbitrários de $\mathfrak{F}$,
  então, a interseção arbitrária $\bigcap\limits_{\alpha \in \Gamma}F_{\alpha}$ pertence a $\mathfrak{F}.$
\end{description}
\end{teo}

\section{Como criar uma Prova}

 \begin{prova}~

 \item i) $\emptyset$ e $X$ pertencem a $\mathfrak{F}$, pois seus
   complementares pertencem a $\tau$.
 \item ii) Se $F_{i}$ é fechado para $i = 1,...,n$. Utilizando a Lei de Morgan, temos:
   \begin{center}
   $(\bigcup\limits_{i=1}^{n} F_{i})^{c} = \bigcap\limits_{i=1}^{n} (F_{i})^{c}$
   \end{center}
   
   Como a interseção $(F_{i})^{c}$ é  um aberto, pois, a interseção
   de conjuntos abertos é aberta, portanto, $\bigcup\limits_{i=1}^{n} F_{i}$ é fechada
   em $X.$
  \item iii) Seja $\{F_{\alpha \in \mathfrak{F} / \alpha \in
   \Gamma}\}$ uma família arbitrária de conjuntos fechados.

  Utilizando a Lei de Morgan, temos:
   \begin{center}
   $(\bigcap_{\alpha \in \Gamma} F_{\alpha})_{c} = \bigcup\limits_{\alpha \in
   \Gamma} F_{\alpha}^{c}$
   \end{center}
   
   Sabemos que $(F_{\alpha})^{c}$ é aberto, assim, a união de
   conjuntos abertos é aberta, então, podemos concluir que $\bigcap\limits_{\alpha \in
   \Gamma}F_{\alpha}$ é fechado em $X$.
 \end{prova}
